The distinction between existence of an index and a uniform index is
already visible in numerical-dimension-zero abundance. Gongyo proved that
a numerically trivial log canonical divisor on a projective slc pair is
rationally linearly trivial \citep[Theorem~1.5]{Gongyo2013}; the resulting
trivializing multiple may depend on the pair. The index conjecture asks
for a common multiple after fixing the dimension and the boundary
coefficients. Its finite-rational-coefficient slc formulation is stated,
for example, in \citet[Conjecture~1.5]{JiangLiu}. The nonnormal setting is
important because slc pairs occur as limits in moduli problems; see
\citet{Kollar2013} and the boundedness theory for polarized slc
Calabi--Yau pairs in \citet{Birkar2023}.

Pluricanonical representations have long connected index problems to
descent. Fujino developed admissible sections in his abundance theorem
for slc threefolds \citep{Fujino2000}, and formulated uniform boundedness
of canonical representations for canonical varieties with trivial canonical
divisor \citep[Conjecture~3.2]{Fujino2001}. His local-index results and
Ishii's global-index theorem for non-klt log canonical threefold pairs
with standard coefficients made this connection explicit
\citep{Fujino2001,Ishii2000}. Here the standard coefficients are
$1$ and $1-1/q$ for positive integers $q$.
Fujino--Gongyo subsequently proved finiteness of the log pluricanonical
representation for every projective lc pair with semiample log canonical
divisor \citep[Theorem~1.1]{FujinoGongyo2014}. Uniform boundedness of these
finite images is the additional issue in the index problem.

Xu organized reductions among the normal index conjecture, its slc
version, and boundedness of pluricanonical representations
\citep{Xu2019}. In particular, his Theorem~1.11 gives the implication
from normal indices in dimension $d$ and bounded representations in
dimension $d-1$ to slc indices in dimension $d$.
Jiang proved the global index bound for boundary-free klt threefolds
\citep[Corollary~1.7]{Jiang2021}.
Jiang--Liu bounded the representations of lc surfaces with a fixed
trivializing multiple, obtaining the slc index theorem in dimensions at
most three and the non-klt lc case in dimension four
\citep[Theorem~1.4 and Corollaries~1.6--1.7]{JiangLiu}.
Xu also treated klt fourfold pairs with nonzero boundary
\citep[Theorem~1.14]{Xu2019}.

Several more recent results give bounds under different geometric
hypotheses. Birkar proved the index conjecture for rationally connected
klt pairs in every dimension, allowing a fixed rational DCC coefficient
set \citep[Corollary~1.8]{Birkar2025}.
Filipazzi--Mauri--Moraga obtained sharp bounds for coregularity-zero
pairs, including slc pairs \citep{FMM2025}; here coregularity zero means
that a dlt model of the normalization with its induced boundary has a
zero-dimensional minimal lc stratum.
Figueroa--Filipazzi--Moraga--Peng developed corresponding index
and complement results in low coregularity \citep{FFMP2025}.
Masamura proved boundedness of the indices of smooth Calabi--Yau
fourfolds \citep[Corollary~1.6]{Masamura2025}.
Uniform bounds need not be numerically small: Singh constructed smooth
Calabi--Yau varieties whose indices grow doubly exponentially with
dimension \citep{Singh2026}.

The proof here follows the established normal-to-slc strategy. It uses
the uniform normal log canonical index theorem of the companion
\emph{Uniform Pluricanonical Iitaka Fibrations}
\citep[Theorem~1.2]{UPI}, together with the all-dimensional good-model
theorem of \emph{Log abundance in characteristic zero}
\citep[Theorem~11.1]{LA}. These are substantial inputs.
Integral cohomology already bounds pluricanonical characters in terms of
middle Betti numbers of suitable covers \citep[Lemma~2.7]{JiangLiu}.
The additional uniformity comes from a rank bound for the rational Hodge
structure generated by a holomorphic top form and the resulting character
bound for arbitrary integral klt pairs with fixed index. The latter is
the klt case of the bounded-representation problem formulated in
\citet[Conjecture~1.1]{JiangLiu}. The final comparison of conductor
residues uses Koll{\'a}r's linking theorem and its even-degree residue
compatibility \citep[Theorem~10 and Proposition~14]{KollarSources}.

The rank argument combines local positivity, moving subvarieties, and
section counts. Its order propagation belongs to the
Ein--K{\"u}chle--Lazarsfeld method \citep{EKL1995}; we use the precise
chain and scalar constructions of \citet[Sections~5--6]{U4} and
\citet[Section~6]{UPI}. Cotangent semipositivity
\citep{GKP2016,CP2019}, divisorial Zariski decomposition
\citep{Nakayama2004,Boucksom2004,BFJ2009}, and effective birationality
\citep{Birkar2023} supply the positivity estimates.
The character argument builds on Matsumura--Wang's decomposition of klt
log Calabi--Yau pairs \citep{MatsumuraWang2025} and the singular
Beauville--Bogomolov theory
\citep{GKP2016,Druel2018,GGK2019,HoringPeternell2019}.
For rationally connected factors we use Han--Jiang's boundedness theorem
and Jiang--Liu's bounded-family representation theorem
\citep{HanJiang,JiangLiu}. The proofs below isolate the new uniform
estimates and explain the adaptations of these earlier constructions.
